\section{\tool{} Design}\label{sec:design}

\tool{} is a static checker for the triple (circuit, generated verifier,
consumer). Its design goal is asymmetric: a \VIOL{} is a warning that later
validation may confirm or refute, but a \VWS{} is a claim that a named binding
property holds for the matched artifacts. We therefore attach to each positive
verdict an explicit evidence requirement, and make every other outcome \UNK{}
with a reason code.

\subsection{Pipeline}

Fig.~\ref{fig:pipeline} shows the pipeline. \emph{Inventory and matching}
enumerates Circom sources, compiled R1CS and symbols, Solidity verifiers and
consumers, and forms candidate triples by matching consumer verifier calls to
verifier contracts and verifiers to circuits by arity and key. Ambiguous or
missing matches are recorded, not guessed. \emph{Extraction} recovers, with
provenance to a file and line: from the circuit, the ordered public signals and
the R1CS with every coefficient and the prime; from the verifier, the arity,
the key constants, and each \code{checkField} call with its exact word offset
and the modulus its helper rejects against; from the consumer, the public
vector as constructed at the call site, the protected effects (transfers,
mints, \code{\_safeMint}, value calls, state writes), replay mappings, and every
guard with its position relative to each effect. A restricted frontend resolves
local modifiers with a single unconditional placeholder and a pure precondition
prefix, unique local base contracts, and exact OpenZeppelin imports; anything
else stays unresolved.%
% src: docs/verdict-rules.md#p1--protected-effect-binding (modifier/inheritance subset); reports/soundness/REPORT-2.md#what-changed
\
\emph{Recognizers} P1--P5 run unconditionally on every retained triple, and
\emph{aggregation} reports per pattern: any applicable \VIOL{} dominates; any
applicable \UNK{} blocks \VWS{}; \VWS{} requires at least one applicable result
and all applicable results verified.%
% src: docs/verdict-rules.md#aggregation-and-non-vacuity-boundary; reports/groundtruth/CLAIMS.md#A3

\begin{figure}[t]
\centering
% F2: BindCheck pipeline. Source: docs/verdict-rules.md; reports/soundness/REPORT-2.md#what-changed;
% dashed validation box: reports/impact/REPORT.md (separate local harness).
\begin{tikzpicture}[
  font=\scriptsize,
  box/.style={draw, rounded corners=2pt, align=center, inner sep=2.5pt, minimum height=0.55cm, fill=white},
  arr/.style={-{Stealth[length=4pt]}, semithick},
]
\node[box, text width=1.6cm] (c) {Circom source\\ + R1CS, symbols};
\node[box, text width=1.6cm, right=0.25cm of c] (v) {snarkjs\\ verifier (.sol)};
\node[box, text width=1.6cm, right=0.25cm of v] (k) {consumer\\ (.sol)};
\node[box, text width=1.2cm, right=0.25cm of k] (pol) {policy flags\\ (explicit)};
\node[box, text width=7.6cm, below=0.35cm of v.south east, anchor=north] (m) {\textbf{Inventory \& triple matching}: verifier calls $\to$ verifier $\to$ circuit (arity, key); ambiguous or missing $\Rightarrow$ recorded};
\node[box, text width=2.3cm, below=0.35cm of m.south west, anchor=north west] (ec) {\textbf{circuit}\\ public signals (ordered);\\ R1CS terms, prime};
\node[box, text width=2.3cm, right=0.12cm of ec] (ev) {\textbf{verifier}\\ arity; $\gamma,\delta,\mathit{IC}$;\\ field guard: offset, modulus};
\node[box, text width=2.3cm, right=0.12cm of ev] (ek) {\textbf{consumer}\\ public vector; effects;\\ replay maps; guards};
\node[box, text width=7.6cm, below=0.35cm of ev.south, anchor=north] (r) {\textbf{Recognizers P1--P5} (all run) $\to$ per triple: \VIOL{} / \VWS{} only with Table~\ref{tab:evidence} evidence / else \UNK{} + reason};
\node[box, text width=7.6cm, below=0.25cm of r] (a) {\textbf{Aggregation} per pattern: \VIOL{} dominates; any applicable \UNK{} blocks \VWS{}};
\node[box, dashed, text width=7.6cm, below=0.45cm of a] (val) {\textbf{Validation} (separate): compile, key correspondence, real proof (E3), unchanged consumer on local chain, committed state (E4), repair + legitimate run};
\draw[arr] (c.south) -- (c.south |- m.north);
\draw[arr] (v.south) -- (v.south |- m.north);
\draw[arr] (k.south) -- (k.south |- m.north);
\draw[arr] (pol.east) -- ++(0.12,0) |- (r.east);
\draw[arr] (ec.north |- m.south) -- (ec.north);
\draw[arr] (ev.north |- m.south) -- (ev.north);
\draw[arr] (ek.north |- m.south) -- (ek.north);
\draw[arr] (ec.south) -- (ec.south |- r.north);
\draw[arr] (ev.south) -- (ev.south |- r.north);
\draw[arr] (ek.south) -- (ek.south |- r.north);
\draw[arr] (r) -- (a);
\draw[arr] (a) -- node[right, font=\tiny, inner sep=1pt]{E1 warnings} (val);
\end{tikzpicture}

\caption{\tool{} pipeline. Solid boxes show the static analyzer. The dashed box tests whether a candidate has proof and execution evidence on a local chain.}
\label{fig:pipeline}
\end{figure}

\subsection{Evidence per verdict}\label{sec:design-evidence}

Table~\ref{tab:evidence} lists the evidence each positive verdict requires.
Three principles recur.
\emph{Path sensitivity:} an equality between an effect value and a public input
counts only if it is enforced on every committing path before the effect---a
preceding \code{require}/\code{assert}, a rejecting inequality branch, or the
effect nested in the equality branch. A guard in a sibling branch does not
count; a guard after the effect is \UNK{} because rollback reasoning is out of
scope.
\emph{Position exactness:} a canonicality guard counts only for the exact
public-vector position from which the replay key derives, only if its helper
rejects against the exact BN254 scalar-field order, and only if the call lies on
the verifier path. A guard on position $i$ is no evidence for a key at
position $j$; a guard against the base field $q$ is no evidence at all.
\emph{Policy is explicit:} \code{msg.sender} binds a recipient only under an
analyst-supplied \code{any\_caller\_may\_receive} policy flag, recorded in the
result.%
% src: docs/verdict-rules.md (all P1-P5 sections)
\
Commitment evidence from R1CS connectivity is recorded but is screen-only: it
never yields \VWS{}, because connectivity does not establish that a
constraint implements Poseidon, nor that Solidity and Circom hash the same
preimage in the same order.

\begin{table}[t]
\centering
\caption{Evidence required for each positive verdict. Everything else is
\UNK{} with a reason code.}
\label{tab:evidence}
\footnotesize
\setlength{\tabcolsep}{2pt}
\begin{tabular}{@{}p{0.55cm}p{\dimexpr\columnwidth-0.55cm-2\tabcolsep\relax}@{}}
\toprule
& \textbf{\VWS{} requires} \\
\midrule
P1 & every recovered recipient, amount, state identity is a verifier public
input, or equated to one by a guard dominating the effect; or recipient is
\code{address(this)}; or \code{msg.sender} under explicit any-caller policy.
Matched circuit and verifier artifacts. Commitments: never. \\
P2 & chain ID, consumer address or domain occurs in the public vector, inside
a hash passed as a public input, or in a guard dominating every effect. \\
P3 & each replay key is proof-bound, checked, consumed, and canonical:
exact-position \code{checkField} against $r$ on the verifier path, a
position-preserving wrapper, or a prior unconditional \code{require(key<r)}. \\
P4 & parsed $\gamma\neq\delta$ pattern, no duplicate $\mathit{IC}$, every
public position guarded against $r$. Parsed key and verification function. \\
P5 & consumer count $=$ verifier arity $=$ circuit public count, and named
signals agree in order; or an exact length guard for dynamic vectors. \\
\bottomrule
\end{tabular}
% src: docs/verdict-rules.md#p1 ... #p5 (rule properties and justifications)
\end{table}

\subsection{Correspondence and extent}\label{sec:design-failures}
The candidate analyzer (content hash \code{753c6e2a\ldots}) compares the
declared verifier receiver with the candidate verifier. A conflict blocks
effect certification and any verifier certification attached to that triple.
An inherited guard supplies evidence only on a resolved forwarding path
that passes the fixed vector unchanged and checks each position. For an
interface check the analyzer must recover the flattened circuit order and
public count. A guarded dynamic vector can establish length compatibility
alone. These checks address observed false certifications. The historical
DEV results in \S\ref{sec:res-dev} include an earlier key-constant error
corrected at the subsequent freeze.
% src: reports/F1/REPORT.md; docs/FREEZE-2.md

For extent binding the static rule covers the case in which a commitment
covers a fixed array but the consumer settles a prefix. It traces the loop
bound and checks whether a dominating guard makes the bound equal to the
committed extent or ensures each skipped row is neutral for the protected
effect. This is source-level evidence under the reported correspondence
assumption. In the Aztec reconstruction the rule distinguishes the faulty
and repaired loops. In the deployed source the verifier input is assembled
in inline assembly and the rule returns \UNK{}. Reliance on unconstrained
elements in the opposite direction remains a separate producer obligation.
% src: reports/S2c/REPORT.md; reports/F1/REPORT.md

\subsection{Local execution oracle}\label{sec:oracle}
An analyst supplies an entitlement policy, a pinned deployment manifest and
an honest reference call. The policy identifies permitted effects and
required obligations independently of the expected verdict. On an owned
local chain the oracle varies the caller, an independent argument, the
encoding of a public word, repeated use, or the deployment context, as the
manifest permits. Each trial records which change was made and compares the
result with the reference under that policy.
% src: oracle/README.md (operators and confirmation rule)

E4 requires a successful transaction whose trace shows the unchanged verifier
returning true on that path. For protected token effects the oracle checks
the transfer record against the recipient's balance change or resulting
token ownership. An event alone or a separate verification of the honest
proof leaves the trial \UNK{}. For coverage, the honest reference must
discharge the required obligation. The oracle then compares the quantity,
token identity and multiplicity discharged when the trial grants the benefit.
For accounting contracts it can compare changes in declared credit and
backing views across the transaction. A reverted or policy-permitted trial
does not confirm a violation. Missing proof, trace or state evidence leaves
the affected question open.
% src: oracle/README.md; reports/F2/REPORT.md
